\chapter{考前 Checklist 与增量路线}
\label{ch:checklists}

本章不是新的理论章，而是把第 \ref{ch:syllabus} 章的官方考纲、第 \ref{ch:primer} 章的理论闭包，以及第 \ref{ch:2025spring}--\ref{ch:2026spring} 章的公开真题解答压缩成一份考前执行手册。读法很简单：先用覆盖表确认没有漏考点，再用四门课的闭卷 checklist 做自测，最后按增量路线把错题回写到理论章或真题章。

\source{路线依据：求真书院 AI 方向博士生资格考试大纲 \parencite{qzc-ai-syllabus-2026}，公开卷 \cite{qzc-ai-2025-spring,qzc-ai-2025-fall,qzc-ai-2026-spring}，以及第 \ref{ch:primer} 章的理论预备。}

本章的蓝色概念词可直接跳回第 \ref{ch:primer} 章的定义、推导或比较表。若某个 checklist 条目只能链接到概念名而没有推导背景，应按 C0 标签补第三章，而不是把背景塞在第七章里。

\section{使用原则}

考前复习最容易出现两种失误：一是只背公式但不会说明公式从哪里来；二是只看真题答案但没有形成可迁移模板。本章采用三个判定标准。

\begin{enumerate}
  \item 能闭卷写出对象、假设、公式和推导骨架。
  \item 能把每个考点映射到至少一道公开卷题或一个可能题型。
  \item 能在错题后判断应该补第 \ref{ch:primer} 章、真题章，还是只补计算熟练度。
\end{enumerate}

每个条目按三档自评：\(0\) 表示只认得名词；\(1\) 表示能复述定义或公式；\(2\) 表示能在新题中独立完成推导或设计。考前最后一轮只追 \(0\to1\) 和高频 \(1\to2\)，不要把所有低频知识点平均用力。

\section{考纲覆盖总表}

下面的表用于确认第七章的路线覆盖，而不是替代第 \ref{ch:primer} 章的教学内容。若某一行自评低于 \(1\)，先回读第 \ref{ch:primer} 章对应小节；若自评为 \(1\) 但真题仍错，再回到第 \ref{ch:2025spring}--\ref{ch:2026spring} 章重做相近题。

\begingroup
\small
\begin{longtable}{L{0.29\textwidth}L{0.30\textwidth}L{0.31\textwidth}}
\toprule
考纲方向 & 必须闭卷掌握 & 回查位置 \\
\midrule
\endfirsthead
\toprule
考纲方向 & 必须闭卷掌握 & 回查位置 \\
\midrule
\endhead
机器学习理论 & \conceptref{concept:pac}{PAC}/\conceptref{concept:agnostic-pac}{agnostic PAC}、\conceptref{concept:vc-sauer}{VC/Sauer}、\conceptref{concept:bias-variance}{模型选择}、线性/逻辑回归、树、最近邻、\conceptref{concept:svm-kernel}{SVM/核}、\conceptref{concept:compressed-sensing}{压缩感知} & 第 \ref{ch:primer} 章学习理论和经典模型；2025 春/秋与 2026 春 A 卷 \\
深度学习和强化学习 & MLP/CNN/Transformer、\conceptref{concept:backprop}{反向传播}、逼近和泛化、视觉任务、无监督学习、\conceptref{concept:generative-models}{生成模型}、\conceptref{concept:bellman}{Bellman}/\conceptref{concept:q-learning}{Q-learning}/\conceptref{concept:policy-gradient}{policy gradient} & 第 \ref{ch:primer} 章深度学习、生成模型、RL；三套 B 卷 \\
人工智能中的优化方法 & \conceptref{concept:convex-first-order}{凸集/凸函数}、次微分、GD/\conceptref{concept:sgd}{SGD}/随机坐标下降、动量、自适应学习率、\conceptref{concept:prox}{prox 递推} & 第 \ref{ch:primer} 章优化；三套 C 卷 \\
自然语言处理 & \conceptref{concept:ngram}{统计建模}、\conceptref{concept:word2vec}{Word2vec}、\conceptref{concept:self-attention}{self-attention}、\conceptref{concept:pretraining-alignment}{预训练对齐}、\conceptref{concept:scaling-law}{scaling law}、\conceptref{concept:rag-mln}{embedding 知识库推理} & 第 \ref{ch:primer} 章 NLP/LLM；三套 D 卷 \\
\bottomrule
\end{longtable}
\endgroup

\section{机器学习理论}

\source{本节对应考纲 Machine Learning Theory。基础理论回查 \textcite{shalev2014uml,mohri2018foundations}，压缩感知回查 \textcite{candes2006compressive}。}

\paragraph{闭卷自测}

\begin{itemize}
  \item 能把固定 \(h\) 的 \(0\)-\(1\) 误差写成 Bernoulli 平均，并快速得到均值、方差、Hoeffding/Chernoff 型界，回查 \conceptref{concept:pac}{PAC 集中界推导}。
  \item 能区分 realizable PAC 与 \conceptref{concept:agnostic-pac}{agnostic PAC}：前者寻找零训练误差假设，后者与 \(\inf_{h\in\Hcal}L_{\Dcal}(h)\) 竞争。
  \item 能从固定假设的集中不等式过渡到有限类 union bound，再过渡到 \conceptref{concept:vc-sauer}{VC/Sauer} 的无限类控制。
  \item 能定义 shattering、增长函数 \(\tau_{\Hcal}(m)\)、\conceptref{concept:vc-sauer}{VC 维}，并用 Sauer 引理解释样本复杂度和 VC 维的关系。
  \item 能求一维几何类的 VC 维：先构造 shatter 下界，再找一个不可实现 pattern 做上界。
  \item 能解释 \conceptref{concept:bias-variance}{模型选择中的 approximation error、estimation error、bias squared 和 variance}。
  \item 能推线性回归正规方程、岭回归一阶条件、逻辑回归负对数似然和梯度。
  \item 能说明决策树的信息增益或 Gini 指标为什么会过拟合，以及剪枝或验证集如何控制复杂度。
  \item 能比较 \(1\)-NN 与 \(k\)-NN 的 bias-variance 和距离尺度敏感性。
  \item 能写 \conceptref{concept:svm-kernel}{hard/soft margin SVM} 的 primal、hinge loss、对偶变量直觉和 kernel PSD 条件。
  \item 能解释 \conceptref{concept:compressed-sensing}{压缩感知} 中的稀疏性、欠定测量、\(\ell_1\) 松弛和 RIP 为什么相关。
\end{itemize}

\paragraph{推导模板}

学习理论题通常不是问“记不记得结论”，而是问能否把概率界从单个对象扩展到整个类。最常用模板是：

\[
  \text{固定 } h
  \Longrightarrow
  \text{集中不等式}
  \Longrightarrow
  \text{union bound}
  \Longrightarrow
  \text{增长函数或复杂度替换}.
\]

其中 \(h\in\Hcal\) 是假设，\(\Hcal\) 是假设类，\(m\) 是样本量。若题目出现 agnostic PAC，还要补上 ERM 三行证明：

\[
  L_{\Dcal}(h_S)
  \le
  L_S(h_S)+\varepsilon
  \le
  L_S(h^*)+\varepsilon
  \le
  L_{\Dcal}(h^*)+2\varepsilon.
\]

这里 \(h_S\) 是 ERM 输出，\(h^*\in\argmin_{h\in\Hcal}L_{\Dcal}(h)\)。第一步和第三步来自统一收敛，第二步来自 ERM 定义。

\paragraph{常见失分点}

\begin{itemize}
  \item 只写 Hoeffding，不写对 \(\Hcal\) 的 \conceptref{concept:pac}{union bound}。
  \item 把 \conceptref{concept:vc-sauer}{VC 维}和参数个数等同，而不说明 shattering。
  \item 忘记 \conceptref{concept:agnostic-pac}{agnostic 情形}要和类内最优竞争，而不是和 Bayes rule 直接竞争。
  \item \conceptref{concept:svm-kernel}{SVM 对偶}中只写 kernel trick，不检查 kernel 是否 symmetric positive semidefinite。
  \item \conceptref{concept:compressed-sensing}{压缩感知}题只背“少量测量可恢复”，不说明稀疏性和 RIP 是恢复保证的条件。
\end{itemize}

\section{深度学习与强化学习}

\source{本节对应考纲 Advanced Deep Learning。神经网络和生成模型回查 \textcite{goodfellow2016deep,bishop2023deep,bishop2006prml}，视觉任务回查 \textcite{szeliski2022vision}，强化学习回查 \textcite{sutton2018rl}。}

\paragraph{闭卷自测}

\begin{itemize}
  \item 能写 MLP、CNN、\conceptref{concept:self-attention}{Transformer} 的输入、输出、参数量和损失函数。
  \item 能推导 \conceptref{concept:backprop}{backprop} 的链式法则形式，并说明复杂度为何与 forward 同阶。
  \item 能解释 sigmoid 饱和、ReLU 死亡、batch normalization 训练/测试统计量差异、residual connection 的梯度通路。
  \item 能区分逼近能力、优化可达性和泛化能力；不会把 universal approximation 当成泛化定理。
  \item 能按输出对象区分特征提取、图像恢复、三维重构、光流估计、识别和分割。
  \item 能写 autoencoder 重构目标和 contrastive learning/InfoNCE 的正负样本结构。
  \item 能说明隐式正则化、迁移学习和元学习分别回答什么问题。
  \item 能写 \conceptref{concept:flow-likelihood}{normalizing flow 的 change-of-variables log-likelihood}。
  \item 能写 \conceptref{concept:vae-elbo}{VAE ELBO}、GAN minimax 目标、DDPM 或 \conceptref{concept:diffusion-score}{score-based diffusion} 的核心训练信号。
  \item 能写 \conceptref{concept:mdp}{MDP 五元组}、\conceptref{concept:bellman}{Bellman 方程}、\conceptref{concept:q-learning}{Q-learning 更新}和 \conceptref{concept:policy-gradient}{policy gradient theorem}。
  \item 能说明 baseline 为什么不改变 \conceptref{concept:policy-gradient}{policy gradient} 的期望。
\end{itemize}

\paragraph{生成模型比较模板}

生成模型比较题优先按数学对象比较，而不是按“效果好不好”比较。

\begingroup
\small
\setlength{\tabcolsep}{3pt}
\begin{longtable}{L{0.14\textwidth}L{0.25\textwidth}L{0.18\textwidth}L{0.25\textwidth}}
\toprule
模型 & 优化对象 & 优势 & 主要风险 \\
\midrule
\endfirsthead
\toprule
模型 & 优化对象 & 优势 & 主要风险 \\
\midrule
\endhead
Flow & \conceptref{concept:flow-likelihood}{显式 likelihood 和 Jacobian determinant} & 可精确求密度 & 可逆结构限制表达形式 \\
VAE & \conceptref{concept:vae-elbo}{ELBO、近似后验和先验匹配} & 概率图清晰，可采样 & posterior collapse、样本模糊 \\
GAN & 生成器和判别器的 minimax game & 样本锐利 & 训练不稳定、mode collapse、无显式 likelihood \\
Diffusion & \conceptref{concept:diffusion-score}{多步去噪或 score matching} & 训练稳定，质量高 & 采样成本高，依赖噪声日程 \\
\bottomrule
\end{longtable}
\endgroup

\paragraph{RL 推导模板}

\conceptref{concept:mdp}{RL 题}先定义状态 \(s\in\mathcal{S}\)、动作 \(a\in\mathcal{A}\)、策略 \(\pi_\theta(a\mid s)\)、奖励 \(r(s,a)\) 和折扣因子 \(\gamma\in[0,1)\)。若问 value function，先写 \conceptref{concept:bellman}{Bellman 分解}；若问策略梯度，先写 \conceptref{concept:policy-gradient}{trajectory likelihood 的 log-derivative trick}。

\[
  \nabla_\theta J(\theta)
  =
  \E_{\tau\sim\pi_\theta}
  \left[
    R(\tau)
    \sum_t
    \nabla_\theta\log \pi_\theta(a_t\mid s_t)
  \right].
\]

其中 \(\tau\) 是轨迹，\(R(\tau)\) 是轨迹回报。baseline \(b(s_t)\) 可减方差，因为

\[
  \E_{a_t\sim\pi_\theta(\cdot\mid s_t)}
  \left[
    b(s_t)\nabla_\theta\log \pi_\theta(a_t\mid s_t)
  \right]
  =
  b(s_t)\nabla_\theta
  \sum_{a_t}\pi_\theta(a_t\mid s_t)
  =
  0.
\]

\paragraph{常见失分点}

\begin{itemize}
  \item 把 \conceptref{concept:backprop}{backprop} 写成口号，不写 \conceptref{concept:jacobian-adjoint}{局部 Jacobian 或 adjoint 递推}。
  \item 只说 Transformer 能并行，不说 \conceptref{concept:self-attention}{self-attention 的 \(Q,K,V\) 和 \(1/\sqrt{d_k}\)}。
  \item \conceptref{concept:vae-elbo}{VAE/ELBO}、\conceptref{concept:flow-likelihood}{flow likelihood}、\conceptref{concept:diffusion-score}{diffusion score matching} 混用符号。
  \item \conceptref{concept:mdp}{RL 题}忘记说明策略、转移、奖励和折扣因子，导致 \conceptref{concept:bellman}{Bellman 方程}没有定义域。
  \item \conceptref{concept:rlhf-dpo}{DPO/RLHF} 题只背名称，不说明偏好数据、reference policy 或 KL regularization 的作用。
\end{itemize}

\section{优化方法}

\source{本节对应考纲 Optimization Methods for AI。凸优化与一阶方法回查 \textcite{nesterov2018convex,beck2017first}，SGD 与学习理论接口回查 \textcite{shalev2014uml}。}

\paragraph{闭卷自测}

\begin{itemize}
  \item 能写 \conceptref{concept:convex-first-order}{凸集、凸函数}、\conceptref{concept:smooth-strong-convex}{\(L\)-smooth、\(\mu\)-strong convex}、subdifferential、\conceptref{concept:prox}{proximal operator} 的定义。
  \item 能从 \conceptref{concept:convex-first-order}{一阶条件}推出凸函数全局最优性。
  \item 能写 Fermat rule \(0\in\partial f(x^*)\)、sum rule 和仿射复合规则。
  \item 能用 \conceptref{concept:smooth-strong-convex}{descent lemma} 推 GD 的下降不等式。
  \item 能用 \conceptref{concept:smooth-strong-convex}{强凸性和光滑性}得到梯度映射收缩或线性收敛。
  \item 能从 \conceptref{concept:prox}{prox optimality} 推 fixed point。
  \item 能用 \conceptref{concept:prox-nonexpansive}{prox 非扩张}推出 one-step inequality。
  \item 能在条件期望下展开 variance decomposition，识别零均值噪声的交叉项为零。
  \item 能解释常数步长 \conceptref{concept:sgd}{SGD} 的 noise floor，以及 minibatch/importance sampling 如何降低方差或 expected smoothness 常数。
  \item 能写 momentum、Nesterov acceleration、AdaGrad、Adam 的状态变量，并说明它们分别改变了哪一部分更新。
\end{itemize}

\paragraph{prox 题核心链条}

若目标是

\[
  \min_{x\in\R^d}
  F(x)=f(x)+R(x),
\]

其中 \(f\) 光滑、\(R\) 凸但可能不可微，\conceptref{concept:prox}{proximal gradient} 的一步是

\[
  x_{k+1}
  =
  \prox_{\gamma R}(x_k-\gamma g_k),
\]

其中 \(\gamma>0\) 是步长，\(g_k\) 是梯度或随机梯度估计。最优点 \(x^*\) 满足

\[
  x^*
  =
  \prox_{\gamma R}(x^*-\gamma\nabla f(x^*)).
\]

考试推导一般沿着下面的链条，其中 prox 最优性与非扩张的完整证明回查 \conceptref{concept:prox-nonexpansive}{第三章 prox 非扩张推导}。

\[
  \text{prox 最优性}
  \Longrightarrow
  \text{fixed point}
  \Longrightarrow
  \text{非扩张}
  \Longrightarrow
  \text{平方距离递推}
  \Longrightarrow
  \text{收敛率或噪声地板}.
\]

\paragraph{常见失分点}

\begin{itemize}
  \item 把 \conceptref{concept:smooth-strong-convex}{\(L\)-smooth 和 \(\mu\)-strong convex} 的不等式方向写反。
  \item 只写 \conceptref{concept:prox}{prox 定义}，不写对应 \conceptref{concept:prox-nonexpansive}{最优性条件}。
  \item \conceptref{concept:sgd}{随机梯度题}不区分条件期望和全期望。
  \item 证明收敛率时漏掉 \conceptref{concept:smooth-strong-convex}{强凸性}提供的误差下界。
  \item Adam/AdaGrad 题只写名称，不写一阶矩、二阶矩或累积梯度变量。
\end{itemize}

\section{自然语言处理}

\source{本节对应考纲 Natural Language Understanding。统计语言模型、embedding、Transformer、LLM 与 RAG 回查 \textcite{jurafsky2026slp,goldberg2017nnnlp,vaswani2017attention,kaplan2020scaling,hoffmann2022training,lewis2020rag,rafailov2023dpo}。}

\paragraph{闭卷自测}

\begin{itemize}
  \item 能写 \conceptref{concept:ngram}{\(n\)-gram MLE}、Markov 假设、backoff/interpolation smoothing 的基本形式。
  \item 能比较 GloVe、\conceptref{concept:word2vec}{Skip-gram}、Node2Vec、TransE 的数据信号、目标函数和归纳偏置。
  \item 能推 \conceptref{concept:word2vec}{Word2vec skip-gram 或 negative sampling} 目标，并解释 embedding 的语义来源。
  \item 能写 \conceptref{concept:self-attention}{self-attention 的 \(Q,K,V\) 公式}，说明矩阵维度和 \(1/\sqrt{d_k}\) 缩放。
  \item 能解释 causal mask、长上下文 attention mask、时间复杂度和显存复杂度。
  \item 能解释 top-\(p\) 相比 greedy、beam、top-\(k\) 的优势和风险。
  \item 能写 \conceptref{concept:pretraining-alignment}{next-token pretraining、SFT}、\conceptref{concept:rlhf-dpo}{RLHF、DPO} 的目标差异。
  \item 能说明 \conceptref{concept:scaling-law}{scaling law} 的经验幂律形式和 compute-optimal trade-off。
  \item 能写 supervised LDA 或 \conceptref{concept:rag-mln}{MLN 的 joint/log-linear probability}。
  \item 能把 \conceptref{concept:rag-mln}{RAG/memory 系统}拆成写入、检索、重排、生成、引用校验和评估闭环。
\end{itemize}

\paragraph{attention 题最小答案}

给输入 \(X\in\R^{N\times d_{\mathrm{model}}}\)，其中 \(N\) 是 token 数，\(d_{\mathrm{model}}\) 是隐藏维度。完整定义和缩放动机回查 \conceptref{concept:self-attention}{Transformer self-attention}。令

\[
  Q=XW_Q,\qquad
  K=XW_K,\qquad
  V=XW_V.
\]

若 \(Q,K\in\R^{N\times d_k}\)，\(V\in\R^{N\times d_v}\)，则 scaled dot-product attention 为

\[
  \operatorname{Attention}(Q,K,V)
  =
  \softmax\left(\frac{QK^\top}{\sqrt{d_k}}\right)V.
\]

这里 \(QK^\top\in\R^{N\times N}\) 是 token 间相似度矩阵。除以 \(\sqrt{d_k}\) 的 motivation 是让点积方差保持常数量级，避免 softmax 过早饱和 \parencite{vaswani2017attention}。

\paragraph{RAG 系统题模板}

\conceptref{concept:rag-mln}{RAG} 题不要只写“加一个向量数据库”。最小可靠闭环是：

\begin{enumerate}
  \item 写入：清洗文档，分块，记录来源、时间、权限和版本。
  \item 检索：用 embedding 召回候选块，并保留 query rewriting 或 hybrid search 的接口。
  \item 重排：用 cross-encoder、规则或多路信号筛掉噪声块。
  \item 生成：把证据块和任务指令组织成 prompt，限制模型不得编造引用。
  \item 校验：检查答案中的 claim 是否被检索证据支持。
  \item 评估：分别看 recall、faithfulness、answer quality、latency 和 cost。
\end{enumerate}

\paragraph{常见失分点}

\begin{itemize}
  \item \conceptref{concept:ngram}{\(n\)-gram} 题只写频率，不说明平滑解决 unseen events。
  \item \conceptref{concept:word2vec}{Word2vec} 题只说“相似词接近”，不写预测上下文或负采样目标。
  \item \conceptref{concept:self-attention}{attention} 题漏写矩阵维度，导致复杂度解释不可检查。
  \item \conceptref{concept:scaling-law}{Scaling law} 题把经验规律写成严格定理，不说明数据、模型和 compute 的条件。
  \item \conceptref{concept:rag-mln}{RAG} 题没有错误模式：检索失败、检索噪声、上下文污染、引用不忠实、权限泄漏。
\end{itemize}

\section{三轮复习路线}

\paragraph{第一轮：理论闭包}

目标是把第 \ref{ch:primer} 章读成可调用的公式库。每个小节结束时，用一张卡片写下四件事：问题 motivation、核心对象、推导入口、考试提醒。若某张卡片只能写出名词，说明还没有真正掌握。

\begingroup
\small
\begin{longtable}{L{0.18\textwidth}L{0.36\textwidth}L{0.36\textwidth}}
\toprule
阶段 & 任务 & 退出标准 \\
\midrule
\endfirsthead
\toprule
阶段 & 任务 & 退出标准 \\
\midrule
\endhead
T-21 至 T-15 & 通读第 \ref{ch:primer} 章，按四门课建立公式卡片 & 每门课至少能闭卷写出 8 个核心对象和 4 个推导入口 \\
T-14 至 T-8 & 重做三套公开卷，不看答案先写 skeleton & 每道题能定位到第 \ref{ch:primer} 章的具体概念链接或标记为新增 gap \\
T-7 至 T-3 & 只处理错题和 \(0/1\) 自评项 & 每个错题都有原因标签和回写位置 \\
T-2 至 T-1 & 做一轮限时模拟和公式口述 & 不再新增大段理论，只修符号、边界条件和常见失分点 \\
\bottomrule
\end{longtable}
\endgroup

\paragraph{第二轮：真题闭环}

每道真题按以下结构重做：

\begin{enumerate}
  \item 不看答案写题意定位：属于哪门课、哪类题型、调用第 \ref{ch:primer} 章哪个工具。
  \item 写最小解题骨架：对象、假设、目标、关键不等式或目标函数。
  \item 对照本书解答，标出缺失的推导步骤。
  \item 把缺失步骤归类为概念缺口、代数缺口、符号缺口、时间管理缺口或题意误读。
  \item 回写到第 \ref{ch:primer} 章 margin note 或自己的错题账。
\end{enumerate}

\paragraph{第三轮：压缩输出}

最后一轮只训练可输出性。每门课准备一页纸，包含：

\begin{itemize}
  \item 五个最可能考的对象定义；
  \item 三个最可能考的推导模板；
  \item 三个最容易犯的符号错误；
  \item 两个可迁移的系统设计框架；
  \item 一个遇到陌生题时的 fallback 策略。
\end{itemize}

\section{错题账与增量生成路线}

每个错题都要有原因标签。没有标签的错题只会变成“再看一遍答案”，不会形成增量。

\begingroup
\small
\begin{longtable}{L{0.16\textwidth}L{0.36\textwidth}L{0.38\textwidth}}
\toprule
标签 & 含义 & 增量动作 \\
\midrule
\endfirsthead
\toprule
标签 & 含义 & 增量动作 \\
\midrule
\endhead
C0 & 官方考纲内，但第 \ref{ch:primer} 章没有讲清 & 补理论章：motivation、定义、推导、引用、考试提醒 \\
C1 & 理论章有，但真题中不会调用 & 在第七章或错题账加入调用模板和相邻真题链接 \\
C2 & 会公式但推导断裂 & 补中间代数步骤，标明每步使用的定理或假设 \\
C3 & 符号或维度错误 & 更新符号约定，加入维度表或变量定义句 \\
C4 & 新公开卷引入新主题 & 先新增 source，再决定补理论章、真题章还是附录 \\
C5 & 系统设计题输出不完整 & 加入 pipeline、failure modes、evaluation metrics 和 trade-offs \\
\bottomrule
\end{longtable}
\endgroup

增量生成时遵循最小改动原则：

\begin{enumerate}
  \item 新增官方 PDF、网页链接和抽取文本，不覆盖旧卷。
  \item 新增一个年份章节，保持“题意定位 + 解答 + 考试提醒”格式。
  \item 若新题引入新知识点，优先补第 \ref{ch:primer} 章，而不是在解答里重复大段背景。
  \item 若只是常见题型变化，优先补本章 checklist 或错题标签。
  \item 每次修改后编译 PDF，检查页数、目录、引用、日志和 PDF 文本抽取。
  \item 对有歧义题目，保留严格数学结论与考试语境解读，避免把猜测写成定理。
\end{enumerate}

\section{考前四十八小时}

最后两天不适合大规模学习新材料，只做可控修补。

\begin{itemize}
  \item 重新写一遍六个骨架：\conceptref{concept:vc-sauer}{PAC/VC}、\conceptref{concept:prox}{prox}-\conceptref{concept:sgd}{SGD}、\conceptref{concept:diffusion-score}{diffusion}/\conceptref{concept:vae-elbo}{ELBO}、\conceptref{concept:policy-gradient}{policy gradient}、\conceptref{concept:self-attention}{self-attention}、\conceptref{concept:rag-mln}{RAG/MLN}。
  \item 每门课各挑一道错题限时重做，超过十五分钟仍无骨架就回看第 \ref{ch:primer} 章对应小节。
  \item 检查所有常用符号：样本量 \(m\)、维度 \(d\)、序列长度 \(N\)、假设类 \(\Hcal\)、分布 \(\Dcal\)、参数 \(\theta\)。
  \item 练习把“我知道”改写成“我能证明”：每个结论至少说出一个假设和一个推导入口。
  \item 系统题只记模板：输入、模块、目标函数或 scoring、failure modes、evaluation。
\end{itemize}

\begin{checkpoint}
第七章的通过标准不是“看完”，而是能在闭卷状态下把任意 checklist 条目展开成一个可评分答案。若只能说概念名，回第 \ref{ch:primer} 章；若会理论但不会套题，回第 \ref{ch:2025spring}--\ref{ch:2026spring} 章；若同类题反复错，按 C0--C5 标签增量修书。
\end{checkpoint}
